%% Separately authored historical derivative following exact source lines 10353--10674.
%% Necessity: exact restoration supplies aleph, CH, and GCH definitions but does
%% not contain the distinct cardinal-operation, beth, cofinality, and continuum
%% catalogs preserved here.
%% Definition: this file is based on historical slices 1--57, 92--182, and
%% 186--221; the overlapping aleph/CH/GCH block is replaced by an eliminable
%% beth definition. The PCF section supplies definitions, elementary proofs,
%% the maximum theorem, and a singular-cardinal application.
%% Standing: this is separately authored reader material, not inherited source
%% or a theorem-status promotion.
%% Nonredundancy: the exact source remains controlling for aleph/CH/GCH, while
%% these retained slices supply additional cardinal arithmetic and calibrations.
%% Use: parts/p06.tex routes this file once at the original cardinal-arithmetic
%% position, and its register table uses the beth hierarchy defined below.

% BEGIN HISTORICAL SLICE 1-57
%% parts/p06/cardinal-arithmetic.tex --- Phase-4 P6 cardinal arithmetic chapter
%% Expands the cardinal material migrated from v1.1 lines 10559--10654.
%% Depth: D-P7-002/003/004 and N-a-035; apparatus: A-061.
%% Singular-cardinal coverage includes Silver and a sourced PCF introduction.
\chapter{Cardinal Arithmetic}
\label{sec:p06-cardinal-arithmetic}
\chaptermark{Cardinal Arithmetic}

\noindent\textit{Reader note.}  This separately authored chapter extends the
exact aleph/CH/GCH source with cardinal operations, cofinality, regular and
singular behavior, and continuum-function calibrations.  Exact inherited
definitions remain controlling where the topics overlap.

Throughout this chapter, cardinal arithmetic means arithmetic of
ordinal cardinals in an ambient theory containing the Axiom of Choice, unless
a statement explicitly says otherwise.

\begin{hbdefinition}{def:p06-cardinal-arithmetic}{Cardinal sum, product, and exponentiation}
For cardinals \(\kappa,\lambda\), define
\[
\begin{aligned}
  \kappa+\lambda
    &\defeq\abs{(\kappa\times\{0\})\cup(\lambda\times\{1\})},\\
  \kappa\cdot\lambda
    &\defeq\abs{\kappa\times\lambda},\\
  \kappa^\lambda
    &\defeq\abs{{}^{\lambda}\kappa},
\end{aligned}
\]
where \({}^{\lambda}\kappa\) is the set of functions from \(\lambda\) to
\(\kappa\).  For a set-indexed family \(\langle\kappa_i:i\in I\rangle\),
\(\sum_{i\in I}\kappa_i\) is the cardinal of the disjoint union and
\(\prod_{i\in I}\kappa_i\) is the cardinal of the Cartesian product.
These operations act on cardinals, not on arbitrary representatives; their
set constructions distinguish them from ordinal addition, multiplication,
and exponentiation.
\index{cardinal arithmetic}
\index{cardinal addition}
\index{cardinal multiplication}
\index{cardinal exponentiation}
\end{hbdefinition}

\begin{hbtheorem}{thm:p06-cardinal-arithmetic-laws}{cited=Jech2003}{Basic laws of cardinal arithmetic}
Assume the Axiom of Choice.  Cardinal addition and multiplication are
commutative and associative, multiplication distributes over addition, and
for cardinals \(\kappa,\lambda,\mu\) with the displayed expressions defined,
\[
  \kappa^{\lambda+\mu}=\kappa^\lambda\cdot\kappa^\mu,
  \qquad
  (\kappa\cdot\lambda)^\mu=\kappa^\mu\cdot\lambda^\mu,
  \qquad
  (\kappa^\lambda)^\mu=\kappa^{\lambda\cdot\mu}.
\]
If \(\kappa\) and \(\lambda\) are infinite, then
\[
  \kappa+\lambda=\kappa\cdot\lambda=\max(\kappa,\lambda).
\]
Moreover \(2^\kappa=\abs{\mathcal P(\kappa)}\) and
\(\kappa<2^\kappa\) for every cardinal \(\kappa\).
\end{hbtheorem}

% END HISTORICAL SLICE 1-57

% Necessity: the exact source defines alephs, CH, and GCH but does not name the
% iterated power-set cardinal sequence used by the later arithmetic register.
% Definition: \(\beth_0=\aleph_0\), \(\beth_{\alpha+1}=2^{\beth_\alpha}\),
% and \(\beth_\lambda=\sup_{\beta<\lambda}\beth_\beta\) at limit ordinals.
% Standing: Beth hierarchy names this transfinite cardinal sequence.
% Nonredundancy: unlike the aleph hierarchy, it follows iterated exponentiation;
% exact-source aleph, CH, and GCH definitions remain controlling and unrepeated.
% Use: Table~\ref{tab:p06-cardinal-arithmetic} records the beth hierarchy as the
% power-set growth comparator for the later continuum-function results.
\begin{hbdefinition}{def:p06-beth-hierarchy}{Beth hierarchy}
The exact aleph hierarchy, CH, and GCH definitions are in
Section~\ref{sec:cumulative-hierarchy-04}.  Define
\[
  \beth_0\defeq\aleph_0,\qquad
  \beth_{\alpha+1}\defeq2^{\beth_\alpha},\qquad
  \beth_\lambda\defeq\sup_{\beta<\lambda}\beth_\beta
\]
at a limit ordinal \(\lambda\).
\index{beth hierarchy}
\end{hbdefinition}

% BEGIN HISTORICAL SLICE 92-182
\begin{hbdefinition}{def:p06-cofinality-regular-singular}{Cofinality and regular or singular cardinals}
A subset \(C\) of an ordinal \(\alpha\) is \emph{cofinal} in \(\alpha\) iff
for every \(\beta<\alpha\) some \(\gamma\in C\) satisfies
\(\beta\leq\gamma\).  The \emph{cofinality} \(\operatorname{cf}(\alpha)\)
is the least ordinal that is the order type of a cofinal subset of \(\alpha\),
equivalently the least ordinal admitting a cofinal map into \(\alpha\).
An infinite cardinal \(\kappa\) is \emph{regular} iff
\(\operatorname{cf}(\kappa)=\kappa\), and \emph{singular} iff
\(\operatorname{cf}(\kappa)<\kappa\).  Cofinality measures how many earlier
stages are required to approach an ordinal; it is not its cardinality.
\index{cofinal subset}
\index{cofinality}
\index{regular cardinal}
\index{singular cardinal}
\end{hbdefinition}

\begin{hbtheorem}{thm:p06-cofinality-facts}{cited=Jech2003}{Basic cofinality facts}
Assume the Axiom of Choice.  For every nonzero limit ordinal \(\alpha\),
\(\operatorname{cf}(\alpha)\) is an infinite regular cardinal and
\(\operatorname{cf}(\alpha)\leq\abs\alpha\).  Every infinite successor
cardinal is regular.  A cardinal \(\kappa\) is singular exactly when it is the
union of fewer than \(\kappa\) many sets, each of cardinality below
\(\kappa\).
\end{hbtheorem}

\begin{hbtheorem}{thm:p06-konig-cardinal}{cited=Konig1905}{K\"onig's theorem}
Assume the Axiom of Choice.  Let \(I\) be a set and let
\(\langle\kappa_i:i\in I\rangle\) and
\(\langle\lambda_i:i\in I\rangle\) be families of cardinals such that
\(\kappa_i<\lambda_i\) for every \(i\in I\).  Then
\[
  \sum_{i\in I}\kappa_i
  <
  \prod_{i\in I}\lambda_i.
\]
In particular, for every infinite cardinal \(\kappa\),
\[
  \operatorname{cf}(2^\kappa)>\kappa
  \qquad\text{and}\qquad
  \kappa^{\operatorname{cf}(\kappa)}>\kappa.
\]
\end{hbtheorem}

\begingroup\small
\begin{longtable}{@{}p{.22\textwidth}p{.37\textwidth}p{.32\textwidth}@{}}
\caption{Cardinal-arithmetic reference table.}
\label{tab:p06-cardinal-arithmetic}\\
\toprule
Entry & Formula or classification & Hypothesis and use \\
\midrule
\endfirsthead
\caption[]{Cardinal-arithmetic reference table (continued)}\\
\toprule
Entry & Formula or classification & Hypothesis and use \\
\midrule
\endhead
\midrule\multicolumn{3}{r}{\emph{continued on next page}}\\
\endfoot
\endlastfoot
Finite sums and products & Definitions in
Definition~\ref{def:p06-cardinal-arithmetic} & Set constructions; ordinal
cardinal assignment uses the local Choice standing. \\
Infinite collapse & \(\kappa+\lambda=\kappa\cdot\lambda=
\max(\kappa,\lambda)\) for infinite \(\kappa,\lambda\) & Axiom of Choice;
supplies the finite-coding bounds used in Part~1. \\
Exponentiation & \(\kappa^\lambda=\abs{{}^{\lambda}\kappa}\),
\(2^\kappa=\abs{\mathcal P(\kappa)}\), and the three exponent laws above & Axiom
of Choice for ordinal-cardinal assignment; Cantor gives
\(\kappa<2^\kappa\). \\
Aleph hierarchy & \(\aleph_0=\omega\), successor by \(\kappa^+\), limit by
supremum & Enumerates the infinite ordinal cardinals under Choice. \\
Beth hierarchy & \(\beth_0=\aleph_0\), successor by power set, limit by
supremum & Tracks iterated power-set cardinalities. \\
Cofinality & Least order type of a cofinal subset & Separates regular
\(\operatorname{cf}(\kappa)=\kappa\) from singular
\(\operatorname{cf}(\kappa)<\kappa\) cardinals. \\
K\"onig's theorem & \(\sum_i\kappa_i<\prod_i\lambda_i\) when
\(\kappa_i<\lambda_i\) & Forces
\(\operatorname{cf}(2^\kappa)>\kappa\); constrains every continuum function. \\
Hausdorff successor recursion &
\((\kappa^+)^\lambda=\kappa^\lambda\cdot\kappa^+\) for infinite
\(\kappa\) and nonzero \(\lambda\) & Reduces a successor base to its
predecessor; the nonzero-exponent hypothesis is essential. \\
CH and GCH & \(2^{\aleph_0}=\aleph_1\); respectively
\(2^{\aleph_\alpha}=\aleph_{\alpha+1}\) for every \(\alpha\) & Independent
of \(\mathsf{ZFC}\); the forcing chapter records the calibrated status. \\
Easton regular-cardinal freedom & Prescribed values of \(2^\kappa\) on
infinite regular cardinals, subject to monotonicity and the cofinality bound &
Theorem~\ref{thm:p06-easton} states the exact cited hypotheses over a
\(\mathsf{ZFC}+\mathsf{GCH}\) ground; no prescription at singular cardinals is
asserted. \\
\bottomrule
\end{longtable}
\endgroup

% END HISTORICAL SLICE 92-182

% BEGIN HISTORICAL SLICE 186-221
\begin{hbtheorem}{thm:p06-hausdorff-recursion}{cited=Hausdorff1904}{Hausdorff's successor recursion}
Assume the Axiom of Choice.  For every infinite cardinal \(\kappa\) and every
nonzero cardinal \(\lambda\),
\[
  (\kappa^+)^\lambda
  =\kappa^\lambda\cdot\kappa^+
  =\max(\kappa^\lambda,\kappa^+).
\]
Equivalently, for ordinals \(\alpha,\beta\),
\[
  \aleph_{\alpha+1}^{\aleph_\beta}
  =\aleph_\alpha^{\aleph_\beta}\cdot\aleph_{\alpha+1}.
\]
The restriction \(\lambda\neq0\) cannot be dropped: the left side would be
\(1\) while the displayed product would be \(\kappa^+\).  Repeated use of
this identity is the Hausdorff recursion for exponentiation with a successor
aleph base \cite{Jech2003}.
\end{hbtheorem}

\begin{hbtheorem}{thm:p06-easton}{cited=Easton1970}{Easton's theorem on regular cardinals}
Let \(F\) be a definable class function from the infinite regular cardinals to
cardinals such that
\(\kappa<\lambda\) implies \(F(\kappa)\leq F(\lambda)\), and
\(\operatorname{cf}(F(\kappa))>\kappa\) for every regular \(\kappa\).
Starting from a ground model of \(\mathsf{ZFC}+\mathsf{GCH}\) carrying such an
\(F\), Easton's forcing yields a cofinality-preserving generic extension in
which \(2^\kappa=F(\kappa)\) at every infinite regular cardinal \(\kappa\).
Thus monotonicity and K\"onig's cofinality condition are, relative to the
existence of the stated ground model, the only restrictions on the continuum
function at regular cardinals.  The theorem does not prescribe exponentiation
at singular cardinals and is a relative-consistency result, not a theorem that
\(\mathsf{ZFC}\) chooses one such function in the ground universe.
\end{hbtheorem}


\begin{hbtheorem}{thm:p06-silver}{cited=Silver1975}{Silver's theorem on singular cardinals}
Let \(\kappa\) be a singular cardinal of uncountable cofinality. If
\(2^\lambda=\lambda^+\) for every infinite cardinal \(\lambda<\kappa\), then
\(2^\kappa=\kappa^+\). Equivalently, a first failure of GCH
cannot occur at a singular cardinal of uncountable cofinality. The
cofinality hypothesis is essential; the theorem does not settle the
countable-cofinality case and does not supply the later PCF analysis.
\end{hbtheorem}

% END HISTORICAL SLICE 186-221

\section{Possible cofinalities}
\label{sec:p06-pcf}

Possible cofinality theory studies the orders obtained from products of
regular cardinals by identifying functions that agree on a large set of
coordinates.  These orders connect regular-cardinal cofinalities with
arithmetic at singular limits.  We continue to assume the Axiom of Choice.

\begin{hbdefinition}{def:p06-pcf-products}{Products of regular cardinals and cofinal families}
Let \(\mathfrak a\) be a nonempty set of infinite regular cardinals,
regarded as initial ordinals.  Its product space is
\[
 \prod\mathfrak a
   =\{f:\mathfrak a\to\bigcup\mathfrak a:
                         f(\kappa)<\kappa\text{ for every }\kappa\in\mathfrak a\}.
\]
This is a set of functions; its cardinality is the cardinal product
\(\prod_{\kappa\in\mathfrak a}\kappa\).
The pointwise order is
\[
 f\leq g
 \quad\Longleftrightarrow\quad
 \forall\kappa\in\mathfrak a\ (f(\kappa)\leq g(\kappa)).
\]
For a nonempty partially ordered set \((P,\leq)\), a subset \(C\)
is cofinal if every \(p\in P\) satisfies \(p\leq c\) for some \(c\in C\).
Define
\[
 \operatorname{cf}(P,\leq)
   =\min\{|C|:C\subseteq P\text{ is cofinal}\}.
\]
It measures the size of a cofinal family, with no assumption that
the family is linearly ordered.
\index{cofinal family}\index{product of regular cardinals}
\end{hbdefinition}

\begin{hbdefinition}{def:p06-pcf-ultrafilter-order}{Ultrafilter quotients of a product}
A proper filter \(D\) on \(\mathfrak a\) contains \(\mathfrak a\),
omits the empty set, and is closed under supersets and finite
intersections.  It is an ultrafilter if for every
\(b\subseteq\mathfrak a\), exactly one of \(b\) and
\(\mathfrak a\setminus b\) belongs to \(D\).
For \(f,g\in\prod\mathfrak a\), put
\[
\begin{aligned}
 f={}_D g
   &\Longleftrightarrow
       \{\kappa:f(\kappa)=g(\kappa)\}\in D,\\
 f\leq_D g
   &\Longleftrightarrow
       \{\kappa:f(\kappa)\leq g(\kappa)\}\in D,\\
 f<_D g
   &\Longleftrightarrow
       \{\kappa:f(\kappa)<g(\kappa)\}\in D.
\end{aligned}
\]
The reduced product \(\prod\mathfrak a/D\) consists of the
\({}_D\)-equivalence classes \([f]_D\), ordered by
\([f]_D\leq[g]_D\) iff \(f\leq_D g\).
For an ultrafilter this is a linear order, and \(<_D\) gives
its strict comparison.
\index{reduced product}\index{ultrafilter quotient}
\end{hbdefinition}

\begin{hbdefinition}{def:p06-possible-cofinalities}{Possible cofinalities and progressive sets}
Define
\[
 \operatorname{pcf}(\mathfrak a)
   =\left\{\operatorname{cf}\!\left(\prod\mathfrak a/D,\leq\right):
                                  D\text{ an ultrafilter on }\mathfrak a\right\}.
\]
The set \(\mathfrak a\) is progressive if
\[
 |\mathfrak a|<\min\mathfrak a.
\]
Principal ultrafilters are included in the definition.
The cofinalities vary with the ultrafilter on one fixed set of coordinates
\cite[Definition~4.1]{Shelah1992Skeptics}.
\index{possible cofinalities}\index{PCF theory}\index{progressive set}
\end{hbdefinition}

\begin{hblemma}{thm:p06-pcf-elementary}{proved=annex:p6.pcf.1}{Elementary properties of possible cofinalities}
Let \(\mathfrak a\) be a nonempty set of infinite regular cardinals.
For every ultrafilter \(D\) on \(\mathfrak a\), the reduced product
is a linear order without a greatest element.
Its cofinality is an infinite regular cardinal \(\lambda\), and
there is a \(<_D\)-increasing cofinal sequence of length \(\lambda\).
Moreover,
\[
 \mathfrak a\subseteq\operatorname{pcf}(\mathfrak a),
 \qquad
 \operatorname{pcf}(\mathfrak a)=\mathfrak a
       \quad\text{if }\mathfrak a\text{ is finite}.
\]
Progressiveness is not needed for these elementary facts.
\end{hblemma}

\begin{hbtheorem}{thm:p06-pcf-maximum}{cited=Shelah1992Skeptics}{PCF maximum theorem}
If \(\mathfrak a\) is a nonempty progressive set of infinite regular
cardinals, then \(\operatorname{pcf}(\mathfrak a)\) has a largest element,
and
\[
 \operatorname{cf}\!\left(\prod\mathfrak a,\leq\right)
       =\max\operatorname{pcf}(\mathfrak a).
\]
Here the order on the left is the pointwise product order.
In particular, some ultrafilter realizes this maximum
\cite[Main Theorem~4.3]{Shelah1992Skeptics}.
\end{hbtheorem}

The theorem identifies the size of a pointwise cofinal family.
The increasing sequence in Lemma~\ref{thm:p06-pcf-elementary} belongs
to an ultrafilter quotient.  These are different order statements.
For example, take \(\mathfrak a=\{\aleph_0,\aleph_1\}\).
Its product is identified with \(\omega\times\omega_1\).
The whole product is a cofinal family of size \(\aleph_1\), and
projection to \(\omega_1\) proves that no smaller family is cofinal.
But a pointwise nondecreasing sequence of length \(\omega_1\) cannot
be cofinal.  Its first coordinates form a nondecreasing sequence of
natural numbers; the first occurrences of its countably many possible
values lie below one countable ordinal.  Beyond that stage the first
coordinate is constant, and a function with a larger first coordinate
is not dominated.  This illustrates the cofinal-family conclusion
without imposing a cofinal-chain conclusion.

\begin{hbcorollary}{thm:p06-pcf-bounds}{proved=annex:p6.pcf.2}{Bounds at a singular limit of regular cardinals}
For a nonempty progressive set \(\mathfrak a\) of infinite regular
cardinals,
\[
 \sup\mathfrak a\leq\max\operatorname{pcf}(\mathfrak a)
                   \leq\left|\prod\mathfrak a\right|.
\]
If \(\mathfrak a\) has no largest element, then
\[
 \sup\mathfrak a<\max\operatorname{pcf}(\mathfrak a).
\]
In particular, for
\(\mathfrak a=\{\aleph_{n+1}:n<\omega\}\),
\[
 \aleph_{\omega+1}
 \leq\max\operatorname{pcf}(\mathfrak a)
 \leq\aleph_\omega^{\aleph_0}
 \leq2^{\aleph_\omega}.
\]
If \(2^{\aleph_\omega}=\aleph_{\omega+1}\), then
\(\max\operatorname{pcf}(\mathfrak a)=\aleph_{\omega+1}\).
The proof uses the cited maximum theorem.
\end{hbcorollary}

The example begins at \(\aleph_1\): it has
\(|\mathfrak a|=\aleph_0<\aleph_1=\min\mathfrak a\).
Adding \(\aleph_0\) to that countable set would remove the strict
progressiveness inequality.

\begin{hbdefinition}{def:p06-strong-limit}{Strong-limit cardinal}
An infinite cardinal \(\lambda\) is a strong limit if
\[
 2^\kappa<\lambda\qquad\text{for every cardinal }\kappa<\lambda.
\]
This concerns power-set sizes below \(\lambda\); its cofinality is
specified separately.
\index{strong-limit cardinal}
\end{hbdefinition}

\begin{hbtheorem}{thm:p06-shelah-aleph-omega-bound}{cited=Shelah1992Skeptics}{Shelah's bound at \(\aleph_\omega\)}
In \(\mathsf{ZFC}\), if \(\aleph_\omega\) is a strong limit, then
\[
 2^{\aleph_\omega}<\aleph_{\omega_4}.
\]
The ordinal \(\omega_4\) is the fourth uncountable initial ordinal,
of cardinality \(\aleph_4\); it is the index of the aleph on the right.
The hypothesis is equivalently
\(2^{\aleph_n}<\aleph_\omega\) for every \(n<\omega\)
\cite[Theorem~A]{Shelah1992Skeptics}.
\end{hbtheorem}

Here \(\operatorname{cf}(\aleph_\omega)=\omega\), so the displayed
Silver theorem does not determine this power-set cardinal.
Shelah's bound is a further theorem about this countable-cofinality
case.  The source develops the PCF argument behind the bound; the
systematic theory is presented in \cite{Shelah1994Cardinal}.
